\documentclass[a4paper,12pt]{article}
\usepackage[utf8]{inputenc}
\usepackage[T1]{fontenc}
\usepackage[ngerman]{babel}
\usepackage{booktabs}
\usepackage{eurosym}
\usepackage{geometry}
\usepackage{hyperref}
\usepackage{xcolor}
\usepackage{array}
\usepackage{calc}

\geometry{margin=2.5cm}

\hypersetup{
  colorlinks=true,
  linkcolor=blue!60!black,
  urlcolor=blue!60!black,
  pdftitle={Kostenplanung Wetterstation-Set},
  pdfauthor={John van Dijk}
}

\title{%
  \vspace{-1.5cm}
  Kostenplanung: Wetterstation-Set\\[6pt]
  \large Materialliste und Preisübersicht für ein Set
}

\date{Stand: März 2026}

\begin{document}
\maketitle

\section*{Übersicht}
Diese Kostenaufstellung listet die benötigten Bauteile für \textbf{ein Wetterstation-Set} gemäß dem Lernmodul \emph{,,Wetterapp-Test 2.0``} auf. Da der BME280-Umweltsensor verhältnismäßig teuer ist, wird zusätzlich eine kostengünstigere Alternative mit einem DHT11-Sensor betrachtet.

Die Preise wurden bei den folgenden Online-Shops recherchiert:
\begin{itemize}
  \item \href{https://www.roboter-bausatz.de}{roboter-bausatz.de}
  \item \href{https://www.reichelt.de}{reichelt.de}
  \item \href{https://www.conrad.de}{conrad.de}
  \item \href{https://www.galaxus.de}{galaxus.de}
\end{itemize}

\noindent

%------------------------------------------------------------------
\section{Stückliste und Preise (Standard-Variante mit BME280)}
%------------------------------------------------------------------
Der BME280 misst Temperatur, Luftfeuchtigkeit und Luftdruck sehr präzise, schlägt sich aber im Preis deutlich nieder.

\begin{table}[h!]
\centering
\renewcommand{\arraystretch}{1.35}
\begin{tabular}{@{} l c r l @{}}
\toprule
\textbf{Bauteil} & \textbf{Anz.} & \textbf{Preis} & \textbf{Link} \\
\midrule

ESP32 NodeMCU 
  & 1 & 8,89\,\euro{}
  & \href{https://www.roboter-bausatz.de/p/nodemcu-esp32-usb-c-38-pin}{RB} \\

BME280 Umweltsensor 
  & 1 & 5,85\,\euro{}
  & \href{https://www.roboter-bausatz.de/p/digital-bme280-temperatur-sensor-luftdruck-feuchtigkeit-i2c-5v-barometer-arduino?number=RBS15806}{RB} \\


Kleines Solarmodul (z.\,B.\ 5V, 1W)
  & 1 & 2,99\,\euro{}
  & \href{https://www.reichelt.de/de/de/shop/produkt/entwicklerboards_-_solarpanel_0_5_w-254385}{Reichelt} \\

Breadboard 830 Kontakte (MB-102)
  & 1 & 1,49\,\euro{}
  & \href{https://www.roboter-bausatz.de/p/mb102-breadboard-830-kontakte}{RB} \\

Jumper-Kabel Set (65\,Stk.)
  & 1 & 2,40\,\euro{}
  & \href{https://www.reichelt.de/de/de/shop/produkt/entwicklerboards_-_kabel-set_verschiedene_laengen_65er-pack_-282690}{Reichelt} \\
  
USB-C-Datenkabel
  & 1 & 2,09\,\euro{}
  & \href{https://www.reichelt.de/de/de/shop/produkt/sync-_ladekabel_usb-a_-_usb-c_0_5_m_schwarz-215846}{Reichelt} \\

\midrule
\multicolumn{2}{l}{\textbf{Zwischensumme Bauteile}}
  & 23,71\,\euro{} & \\
\multicolumn{2}{l}{Versandkosten abgerundet (geschätzt)}
  & 6,00\,\euro{} & \\
\midrule
\multicolumn{2}{l}{\textbf{Gesamt (inkl. Versand)}}
  & \textbf{29,71\,\euro{}} & \\
\bottomrule
\end{tabular}
\label{tab:kosten-bme280}
\end{table}

\newpage
%------------------------------------------------------------------
\section{Alternative: Kostengünstigere Variante mit DHT11}
%------------------------------------------------------------------
Um Kosten zu sparen, kann der teure BME280 durch einen einfachen DHT11-Sensor ersetzt werden. Der DHT11 misst ebenfalls Temperatur und Luftfeuchtigkeit, erfasst jedoch \textbf{keinen Luftdruck} und ist weniger präzise. Für den grundlegenden Aufbau einer Wetterstation im Unterricht ist dies jedoch oft ausreichend und reduziert die Kosten pro Set erheblich.

\begin{table}[h!]
\centering
\renewcommand{\arraystretch}{1.35}
\begin{tabular}{@{} l c r l @{}}
\toprule
\textbf{Bauteil} & \textbf{Anz.} & \textbf{Preis} & \textbf{Link} \\
\midrule

ESP32 NodeMCU
  & 1 & 8,89\,\euro{}
  & \href{https://www.roboter-bausatz.de/p/nodemcu-esp32-usb-c-38-pin}{RB} \\

\textbf{DHT11 Sensor (Temperatur/Feuchte)}
  & 1 & \textbf{1,85\,\euro{}}
  & \href{https://www.roboter-bausatz.de/p/dht11-digitaler-feuchtigkeits-und-temperatursensor}{RB} \\

Kleines Solarmodul
  & 1 & 2,99\,\euro{}
  & \href{https://www.reichelt.de/de/de/shop/produkt/entwicklerboards_-_solarpanel_0_5_w-254385}{Reichelt} \\

Breadboard 830 Kontakte
  & 1 & 1,49\,\euro{}
  & \href{https://www.roboter-bausatz.de/p/mb102-breadboard-830-kontakte}{RB} \\

Jumper-Kabel Set (65\,Stk.)
  & 1 & 2,40\,\euro{}
  & \href{https://www.reichelt.de/de/de/shop/produkt/entwicklerboards_-_kabel-set_verschiedene_laengen_65er-pack_-282690}{Reichelt} \\

USB-C-Datenkabel
  & 1 & 2,09\,\euro{}
  & \href{https://www.reichelt.de/de/de/shop/produkt/sync-_ladekabel_usb-a_-_usb-c_0_5_m_schwarz-215846}{Reichelt} \\

\midrule
\multicolumn{2}{l}{\textbf{Zwischensumme Bauteile}}
  & 19,71\,\euro{} & \\
\multicolumn{2}{l}{Versandkosten abgerundet (geschätzt)}
  & 6,00\,\euro{} & \\
\midrule
\multicolumn{2}{l}{\textbf{Gesamt (inkl. Versand)}}
  & \textbf{25,71\,\euro{}} & \\
\bottomrule
\end{tabular}
\label{tab:kosten-dht11}
\end{table}

\smallskip\noindent
\textit{Zusätzlicher Hinweis zur Software:} Der Code in der \emph{Wetterapp-Test 2.0} Anleitung müsste für die Nutzung des DHT11-Sensors geringfügig angepasst werden, da auf die Adafruit DHT Bibliothek statt auf die BME280 Bibliothek zurückgegriffen wird.

\section{Preisvergleich -- weitere Shops}
%------------------------------------------------------------------

Die folgende Tabelle zeigt beispielhafte Einzelpreise bei den anderen Anbietern, soweit dort vergleichbare Produkte verfügbar sind.

\begin{table}[h!]
\centering
\renewcommand{\arraystretch}{1.35}
\begin{tabular}{@{} l r r r @{}}
\toprule
\textbf{Bauteil}
  & \textbf{reichelt.de}
  & \textbf{conrad.de}
  & \textbf{galaxus.de} \\
\midrule

ESP32 NodeMCU
  & ab~8,00\,\euro{}
  & ab~10,00\,\euro{}
  & ab~9,00\,\euro{} \\

BME280 Sensor
  & ab~12,00\,\euro{}
  & ab~15,00\,\euro{}
  & ab~12,00\,\euro{} \\

DHT11 Sensor
  & ab~3,00\,\euro{}
  & ab~4,50\,\euro{}
  & ab~3,50\,\euro{} \\

Breadboard 830 Kontakte
  & 3,85\,\euro{}
  & 4,49\,\euro{}
  & -- \\

\bottomrule
\end{tabular}
\label{tab:vergleich}
\end{table}

\smallskip\noindent
\textit{Hinweis: Bei conrad.de und galaxus.de können Markenboards (z.\,B. von Adafruit oder Joy-IT) angeboten werden, wodurch die Preise deutlich höher ausfallen. Für den Unterrichtseinsatz sind die günstigen kompatiblen Modelle von roboter-bausatz.de in der Regel ausreichend.}


%------------------------------------------------------------------
\section{Kalkulationsbeispiel: Klassensatz (15 Sets)}
%------------------------------------------------------------------

Hier vergleichen wir einen Klassensatz mit der BME280- und der DHT11-Variante. Da Lötmaterial und Jumper-Kabel oft im Bestand sind, können die tatsächlichen Kosten variieren. In dieser Aufstellung wird von kompletten Neuanschaffungen für alle 15 Gruppen ausgegangen.

\begin{table}[h!]
\centering
\renewcommand{\arraystretch}{1.35}
\begin{tabular}{@{} l r r @{}}
\toprule
\textbf{Bauteil (15 Sets)} & \textbf{Mit BME280} & \textbf{Mit DHT11} \\
\midrule
ESP32 NodeMCU (15\,x 8,89\,\euro{}) & 133,35\,\euro{} & 133,35\,\euro{} \\
Sensoren (BME280: 5,85\,\euro{} | DHT11: 1,85\,\euro{}) & 87,75\,\euro{} & 27,75\,\euro{} \\
Solarmodule (15\,x 2,99\,\euro{}) & 44,85\,\euro{} & 44,85\,\euro{} \\
Breadboards (15\,x 1,49\,\euro{}) & 22,35\,\euro{} & 22,35\,\euro{} \\
Jumper-Kabel (15\,x 2,40\,\euro{}) & 36,00\,\euro{} & 36,00\,\euro{} \\
USB-C-Datenkabel (15\,x 2,09\,\euro{}) & 31,35\,\euro{} & 31,35\,\euro{} \\
\midrule
\textbf{Zwischensumme Bauteile} & 355,65\,\euro{} & 295,65\,\euro{} \\
Versandkosten (RB + Reichelt) & 9,85\,\euro{} & 9,85\,\euro{} \\
\midrule
\textbf{Gesamtpreis (15 Sets)} & \textbf{365,50\,\euro{}} & \textbf{305,50\,\euro{}} \\
\textit{Preis pro Set} & \textit{$\approx$\,24,37\,\euro{}} & \textit{$\approx$\,20,37\,\euro{}} \\
\bottomrule
\end{tabular}
\caption{Kalkulationsbeispiel für einen Klassensatz (15 Gruppen)}
\label{tab:klassensatz}
\end{table}

%------------------------------------------------------------------
\section{Synergie-Effekte (Kombination mit Fotometer-Set)}
%------------------------------------------------------------------
Wenn die Bauteile aus dem Lernmodul \emph{,,Vom LDR zum Fotometer``} bereits vorhanden sind, können einige Artikel (wie Breadboard und Jumper-Kabel) wiederverwendet werden. Dadurch entfallen Kosten in Höhe von \textbf{3,89\,\euro{}} pro Set (1,49\,\euro{} + 2,40\,\euro{}). 

Die reinen Materialkosten für ein Aufrüst-Set sinken entsprechend:
\begin{itemize}
  \item \textbf{Mit BME280-Sensor:} von 23,71\,\euro{} auf \textbf{19,82\,\euro{}}
  \item \textbf{Mit DHT11-Sensor:} von 19,71\,\euro{} auf \textbf{15,82\,\euro{}}
\end{itemize}

\end{document}
